
\documentclass{article}
\usepackage{colortbl}
\usepackage{makecell}
\usepackage{multirow}
\usepackage{supertabular}

\begin{document}

\newcounter{utterance}

\centering \large Interaction Transcript for game `dond', experiment `coop\_en', episode 10 with reflexion{-}Qwen3.5{-}4B.
\vspace{24pt}

{ \footnotesize  \setcounter{utterance}{1}
\setlength{\tabcolsep}{0pt}
\begin{supertabular}{c@{$\;$}|p{.15\linewidth}@{}p{.15\linewidth}p{.15\linewidth}p{.15\linewidth}p{.15\linewidth}p{.15\linewidth}}
    \# & $\;$A & \multicolumn{4}{c}{Game Master} & $\;\:$B\\
    \hline

    \theutterance \stepcounter{utterance}  
    & & \multicolumn{4}{p{0.6\linewidth}}{
        \cellcolor[rgb]{0.9,0.9,0.9}{
            \makecell[{{p{\linewidth}}}]{
                \texttt{\tiny{[P1$\langle$GM]}}
                \texttt{You have attempted this same task before and did not succeed. These are the plans you wrote after those attempts.} \\
\texttt{Trial \#1: New plan: I failed because I prematurely submitted a secret proposal with only the wires, leaving myself with the chalks (value 0) and the glass (value 8) unallocated, which resulted in a non{-}complementary split or a suboptimal outcome where the total value wasn't maximized. I misunderstood that I needed to explicitly assign the high{-}value glass to the other player in my proposal to ensure the split was complementary and maximized the sum. Next time, I will first analyze the item values (Glass=8, Wires=2, Chalks=0) and immediately submit a secret proposal that assigns the glass to the other player and the wires to myself, ensuring the proposals are complementary and the total value is maximized without needing further negotiation.} \\
\\ 
\\ 
\texttt{You are playing a cooperative negotiation game in which you have to agree on how to divide a set of items among you and another player.} \\
\\ 
\texttt{Rules:} \\
\texttt{(a) You and the other player are given a set of items. Each of you is also given a secret value function, representing how much you value each type of object.} \\
\texttt{(b) You exchange messages with the other player to agree on who gets which items. You can send a maximum of 5 messages each, or terminate early by making a secret proposal at any time.} \\
\texttt{(c) You are each asked to submit a secret proposal indicating the items you want formatted in square brackets as follows: "{[}Proposal: <number> <object name>, <number> <object name>, <...>{]}"} \\
\texttt{(d) If your proposals are complementary, i.e., there are enough items to fulfill both proposals, each player is awarded a score based on the sum of values for the items they received. Otherwise, both of you get zero points.} \\
\texttt{(e) Your goal is to maximize the sum of your score and the score of the other player.} \\
\\ 
\texttt{Let us start.} \\
\\ 
\texttt{The set of available items is:} \\
\\ 
\texttt{1 glass, 2 wires, 2 chalks} \\
\\ 
\texttt{Your secret value function is:} \\
\\ 
\texttt{glass: 8, wire: 1, chalk: 0} \\
\\ 
\texttt{IMPORTANT: Your messages, unless it is the secret proposal, are directly transmitted to the other player, so do not include any response to the rules or text announcing your message. To make a secret proposal, use the indicated format. Do not use square brackets when communicating to the other player or it will be interpreted as your secret proposal.} \\
            }
        }
    }
    & \\ \\

    \theutterance \stepcounter{utterance}  
    & \multicolumn{4}{p{0.6\linewidth}}{
        \cellcolor[rgb]{0.9,0.9,0.9}{
            \makecell[{{p{\linewidth}}}]{
                \texttt{\tiny{[P1$\rangle$GM]}}
                \texttt{I will propose a split that maximizes the total value while ensuring the proposals are complementary.} \\
\\ 
\texttt{My value function: Glass=8, Wire=1, Chalk=0.} \\
\texttt{Total value of all items = 8 + 2*1 + 2*0 = 10.} \\
\\ 
\texttt{To maximize the sum, I should aim for a split where the total value is 10. Since I value the glass highly (8) and the wires slightly (1), and the other player likely values them similarly or at least non{-}negatively, I should offer the glass to the other player to ensure we both get something valuable. If I keep the glass, the other player might get nothing valuable, potentially leading to a low sum if they reject or if the split isn't complementary in a way that satisfies both.} \\
\\ 
\texttt{Let's try this split:} \\
\texttt{{-} I get: 2 wires (Value = 2)} \\
\texttt{{-} Other player gets: 1 glass (Value = 8)} \\
\texttt{{-} Chalks (Value 0) can be given to either or ignored as they add no value.} \\
\\ 
\texttt{This split is complementary:} \\
\texttt{{-} My proposal: {[}2 wires{]}} \\
\texttt{{-} Other player's proposal: {[}1 glass, 2 chalks{]} (or just {[}1 glass{]})} \\
\\ 
\texttt{Wait, the rules say "If your proposals are complementary, i.e., there are enough items to fulfill both proposals". This usually means the union of items in both proposals must equal the set of all available items, or at least not overlap in a way that leaves items unassigned if the goal is to divide the set. Actually, re{-}reading rule (d): "If your proposals are complementary... each player is awarded a score based on the sum of values for the items they received." This implies the items in my proposal + items in the other player's proposal must cover all items (or at least the items I am claiming and they are claiming must not overlap and must sum to the total set).} \\
\\ 
\texttt{Let's refine the plan to be safe and maximize value.} \\
\texttt{I will propose taking the 2 wires.} \\
\texttt{The other player will take the 1 glass and the 2 chalks.} \\
\texttt{Total value for me: 2.} \\
\texttt{Total value for other: 8 + 0 = 8.} \\
\texttt{Total sum: 10.} \\
\\ 
\texttt{If I propose taking the glass (8) and the other player takes the wires (2), the sum is 10 as well. However, since I value the glass more, I might prefer it. But if the other player values the glass even more, they might not accept it. The prompt says "Your goal is to maximize the sum of your score and the score of the other player." This is a cooperative game. So I should propose a split that gives the other player something they value, even if I value it less, to ensure the total is maximized.} \\
\\ 
\texttt{Assuming} \\
            }
        }
    }
    & & \\ \\

    \theutterance \stepcounter{utterance}  
    & & & \multicolumn{2}{p{0.3\linewidth}}{
        \cellcolor[rgb]{0.9,0.9,0.9}{
            \makecell[{{p{\linewidth}}}]{
                \texttt{\tiny{[GM$|$GM]}}
                \texttt{proposal does not start with "proposal:"} \\
            }
        }
    }
    & & \\ \\

    \theutterance \stepcounter{utterance}  
    & \multicolumn{4}{p{0.6\linewidth}}{
        \cellcolor[rgb]{0.95,0.95,0.95}{
            \makecell[{{p{\linewidth}}}]{
                \texttt{\tiny{[P2$\rightarrow$P1]}}
                \texttt{Post{-}episode Reflexion (not sent to GM):} \\
\texttt{New plan: I failed because I sent a conversational explanation and a partial proposal instead of a direct, compliant message. The rules strictly forbid square brackets in regular communication (as they trigger a secret proposal) and require concise, direct statements without meta{-}commentary. Next time, I will immediately submit a secret proposal in the exact format "{[}Proposal: <number> <object name>, ...{]}" without any preceding text, explanation, or square brackets in the message body, ensuring the proposal assigns the high{-}value glass to the other player and the wires to myself to guarantee a complementary split and maximize the total score.} \\
            }
        }
    }
    & & \\ \\

\end{supertabular}
}

\end{document}
